\section{Jacobian 的构造}
\label{sec:jacobian-construction}
% ============================================================================

现在终于来到本章的核心构造。
给定一条亏格为 $g$ 的紧黎曼面 $X$，前面章节已经告诉我们：
$H^0(X,\Omega^1)$ 是一个 $g$ 维复向量空间。
如果只停在这里，我们知道“有多少条全纯微分”，却还没有一个真正的几何对象去承载这些微分的积分信息。

\textbf{我们遇到了什么问题？}
微分本身只是线性对象；
而我们真正想记录的是：沿着闭路径积分以后会得到哪些周期？
这些周期不是随意散落的复数，它们会组成一个格。
如果不把这个格也记下来，我们就得不到一个紧对象，更谈不上后续的 Abel 簇结构与 Hecke 作用。

\textbf{这个概念的核心思想是什么？}
Jacobian 簇的想法非常简单：
先把所有全纯微分放进向量空间 $H^0(X,\Omega^1)^*$，
再把所有闭路径积分产生的周期收集成一个格 $\Lambda_X$，最后取商
\[
  J(X)=H^0(X,\Omega^1)^*/\Lambda_X.
\]
也就是说，Jacobian 就是“全纯微分的积分空间模掉周期不确定性”。

\textbf{引入之后能得到什么？}
这样一来，曲线上的点可以通过积分送进 $J(X)$，这就是 Abel--Jacobi 映射（Abel--Jacobi map）\index{Abel--Jacobi映射!Abel--Jacobi map}。
当 $g=1$ 时，这个商正好恢复原来的椭圆曲线；
当 $X=X_0(N)$ 时，$H^0(X_0(N),\Omega^1)$ 又与 $\Sk_2(\GammaO(N))$ 相同，
于是模形式直接写进了 Jacobian 的定义里。
这就是本章真正的桥梁时刻。

\begin{figure}[ht]
\centering
\begin{tikzpicture}[>=Stealth, font=\small]
  % curve
  \begin{scope}[shift={(0.8,1.8)}]
    \draw[thick, fill=green!8]
      plot[smooth cycle, tension=0.85] coordinates {(-1.2,0.0) (-0.8,1.0) (0.4,1.1) (1.2,0.3) (0.9,-0.9) (-0.1,-1.1) (-1.0,-0.6)};
    \fill (0.95,0.15) circle (1.5pt) node[right] {$P$};
    \fill (-0.65,-0.55) circle (1.5pt) node[left] {$P_0$};
    \draw[->, thick] (-0.55,-0.5) .. controls (0.0,-0.9) and (0.6,-0.6) .. (0.85,0.05);
    \node at (0.1,-1.25) {$X$};
  \end{scope}

  \draw[->, very thick] (2.8,1.8) -- (5.2,1.8) node[midway, above] {$P\mapsto \left(\omega\mapsto \int_{P_0}^{P}\omega\right)$};
  \node at (4.0,2.45) {Abel--Jacobi 映射};

  % lattice / jacobian
  \begin{scope}[shift={(7.8,0.5)}]
    \draw[thick, fill=blue!6] (0,0) rectangle (3.0,2.6);
    \draw[step=1.0, thin, gray!70] (0,0) grid (3.0,2.6);
    \node at (1.5,2.95) {$H^0(X,\Omega^1)^*$};
    \draw[decorate, decoration={brace, amplitude=5pt}, thick]
      (3.2,0) -- (3.2,2.6) node[midway, right=6pt] {模去 $\Lambda_X$};
    \begin{scope}[shift={(4.8,1.3)}]
      \draw[thick, fill=orange!10]
        plot[smooth cycle, tension=0.85] coordinates {(-1.0,0.0) (-0.5,0.65) (0.3,0.7) (0.95,0.25) (1.05,-0.3) (0.4,-0.75) (-0.45,-0.7) (-0.95,-0.2)};
      \node at (0,-1.15) {$J(X)$};
    \end{scope}
  \end{scope}
\end{tikzpicture}
\caption{Jacobian 的解析构造：曲线 $X$ 上从基点 $P_0$ 到点 $P$ 的积分给出一个向量；模去所有闭路径产生的周期格 $\Lambda_X$ 后，就得到点 $\phi_{P_0}(P)\in J(X)$。}
\label{fig:ch06-jacobian}
\end{figure}


\begin{definition}[周期格]
\label{def:period-lattice}
设 $X$ 是亏格为 $g$ 的紧黎曼面，记
\[
  V_X=H^0(X,\Omega^1)^*.
\]
对每个同调类 $\gamma\in H_1(X,\Z)$，定义线性泛函
\[
  I_\gamma: H^0(X,\Omega^1)\longrightarrow \C,
  \qquad
  I_\gamma(\omega)=\int_\gamma \omega.
\]
由此得到群同态
\[
  \iota_X:H_1(X,\Z)\longrightarrow V_X,
  \qquad
  \gamma\longmapsto I_\gamma.
\]
其像
\[
  \Lambda_X=\iota_X\bigl(H_1(X,\Z)\bigr)
\]
称为 $X$ 的\textbf{周期格（period lattice）}\index{周期格!period lattice}。
\end{definition}

\begin{proposition}[周期格确实是一个格]
\label{prop:period-lattice-is-lattice}
设 $X$ 的亏格为 $g$。
则 $\iota_X$ 是单射，并且 $\Lambda_X\subset V_X$ 是秩 $2g$ 的离散子群。
因此 $V_X/\Lambda_X$ 是一个 $g$ 维复环面。
\end{proposition}

\begin{proof}
先证单射。
若 $\gamma\in H_1(X,\Z)$ 满足 $I_\gamma=0$，则对任意全纯微分 $\omega$ 都有
\[
  \int_\gamma \omega=0.
\]
取复共轭可得对任意反全纯微分 $\overline{\omega}$ 也有
\[
  \int_\gamma \overline{\omega}=0.
\]
而紧黎曼面上的调和 $1$-形式恰好分解为“全纯部分 $+$ 反全纯部分”，
所以 $\gamma$ 与每个调和 $1$-形式的积分都为零。
调和形式代表了 $H^1(X,\C)$ 的全部上同调类，
因此 $\gamma$ 与 $H^1(X,\C)$ 中每个类的配对都为零。
由 de Rham 配对
\[
  H_1(X,\C)\times H^1(X,\C)\longrightarrow \C
\]
的非退化性，得到 $\gamma=0$ 在 $H_1(X,\C)$ 中成立。
由于 $H_1(X,\Z)$ 是自由阿贝尔群，故 $\gamma=0$。

再看离散性。
将 $\iota_X$ 张量到实数，得到实线性映射
\[
  \iota_{X,\R}:H_1(X,\R)\longrightarrow V_X.
\]
上面的论证同样说明它仍然单射。
而两边的实维数都等于 $2g$：
左边因为 $H_1(X,\Z)$ 秩为 $2g$，右边因为 $V_X$ 的复维数为 $g$。
故 $\iota_{X,\R}$ 是实线性同构。
于是 $\Lambda_X$ 作为整数格 $H_1(X,\Z)$ 在这个同构下的像，必是 $V_X$ 中的一个秩 $2g$ 离散子群。
所以 $V_X/\Lambda_X$ 是一个 $g$ 维复环面。
\end{proof}

\begin{definition}[Jacobian 簇]
\label{def:jacobian}
设 $X$ 是亏格为 $g$ 的紧黎曼面。
定义它的\textbf{Jacobian 簇（Jacobian variety）}\index{Jacobian簇!Jacobian variety}
为复环面
\[
  J(X)=V_X/\Lambda_X
  =H^0(X,\Omega^1)^*/H_1(X,\Z).
\]
Riemann 双线性关系进一步说明 $J(X)$ 带有自然主极化，
因此它实际上是一个 $g$ 维 Abel 簇。
\end{definition}

\textbf{注。}
Abel 定理告诉我们，$J(X)$ 也可以看成次数为 $0$ 的除子类群 $\operatorname{Pic}^0(X)$。
本节我们主要使用解析构造，因为它最直接地展示了模形式如何进入定义。

\begin{definition}[Abel--Jacobi 映射]
\label{def:abel-jacobi-map}
固定基点 $P_0\in X$。
对每个点 $P\in X$，任选一条从 $P_0$ 到 $P$ 的路径 $\eta$，定义
\[
  \phi_{P_0}(P)(\omega)=\int_\eta \omega
  \qquad \bmod \Lambda_X.
\]
由此得到映射
\[
  \phi_{P_0}:X\longrightarrow J(X),
\]
称为以 $P_0$ 为基点的\textbf{Abel--Jacobi 映射}。
\end{definition}

\begin{proposition}[Abel--Jacobi 映射的基本性质]
\label{prop:abel-jacobi-basic}
映射 $\phi_{P_0}$ 良定义且全纯。
若把 $P_0$ 改成另一个基点 $Q_0$，得到的新映射只与 $\phi_{P_0}$ 相差 $J(X)$ 上的一个平移。
\end{proposition}

\begin{proof}
先证良定义。
若 $\eta_1,\eta_2$ 都从 $P_0$ 连到 $P$，则它们的差 $\eta_1-\eta_2$ 构成闭路径，因而给出某个同调类 $\gamma\in H_1(X,\Z)$。
对任意 $\omega\in H^0(X,\Omega^1)$，有
\[
  \int_{\eta_1}\omega-\int_{\eta_2}\omega
  =\int_{\gamma}\omega=I_\gamma(\omega),
\]
这正是周期格中的元素。
因此两条路径得到的值在 $J(X)=V_X/\Lambda_X$ 中相同。

再证全纯性。
取一点 $P\in X$ 的局部坐标 $z$，并令 $z(P)=0$。
若 $\omega=f(z)\,dz$ 是任一全纯微分，则
\[
  \int_{P_0}^{Q}\omega
  =C+\int_0^{z(Q)} f(\zeta)\,d\zeta
\]
在局部是 $z(Q)$ 的全纯函数。
对一组基微分同时成立，因此 $\phi_{P_0}$ 在局部由若干全纯函数给出，是全纯映射。

若改用基点 $Q_0$，则对任意 $P$ 有
\[
  \phi_{Q_0}(P)-\phi_{P_0}(P)
  =\left(\omega\mapsto \int_{Q_0}^{P}\omega-\int_{P_0}^{P}\omega\right)
  =\left(\omega\mapsto \int_{Q_0}^{P_0}\omega\right),
\]
右边与 $P$ 无关。
故新旧映射只差一个常向量，即 Jacobian 上的平移。
\end{proof}

\begin{theorem}[亏格 $1$ 时曲线等于其 Jacobian]
\label{thm:genus1-jacobian}
设 $X$ 是亏格 $1$ 的紧黎曼面，并固定基点 $P_0\in X$。
则 Abel--Jacobi 映射
\[
  \phi_{P_0}:X\longrightarrow J(X)
\]
是双全纯同构。
特别地，亏格 $1$ 曲线的 Jacobian 就是它自身；
若 $X$ 还带有代数结构与基点，则这正恢复熟悉的椭圆曲线群结构。
\end{theorem}

\begin{proof}
因为 $g=1$，空间 $H^0(X,\Omega^1)$ 一维，取非零微分 $\omega$ 作为基。
其零点除子的次数等于规范除子的次数 $2g-2=0$，
而零点次数总是非负，故 $\omega$ 没有零点。

设 $\widetilde{X}\to X$ 是普遍覆盖。
由于 $\widetilde{X}$ 单连通，函数
\[
  F(Q)=\int_{\widetilde{P}_0}^{Q} \widetilde{\omega}
\]
在 $\widetilde{X}$ 上良定义，其中 $\widetilde{\omega}$ 是 $\omega$ 的拉回。
并且
\[
  dF=\widetilde{\omega}.
\]
因为 $\widetilde{\omega}$ 处处不为零，$F$ 的导数处处非零，故 $F$ 是局部双全纯映射。
其像是 $\C$ 中一个既开又闭的连通子集，只能是整个 $\C$。
因此 $F:\widetilde{X}\to\C$ 是覆盖映射。
但 $\C$ 单连通，故该覆盖只能是双全纯同构。
所以 $\widetilde{X}\cong\C$。

现在看 deck transformation 的作用。
若 $\delta$ 是覆盖变换，则
\[
  F(\delta Q)-F(Q)
\]
的导数为零，因此它是常数；这个常数恰是沿相应闭路径对 $\omega$ 的积分，也就是一个周期。
于是 deck transformation 群恰好通过平移作用，其平移向量组成周期格 $\Lambda_X$。
因此
\[
  X\cong \C/\Lambda_X.
\]
另一方面，由于 $H^0(X,\Omega^1)$ 一维，$V_X=H^0(X,\Omega^1)^*$ 也可与 $\C$ 识别；
在此识别下 Jacobian 恰为
\[
  J(X)=\C/\Lambda_X.
\]
由构造可知 $\phi_{P_0}$ 正是这个商同构。
故它是双全纯同构。
\end{proof}

\begin{proposition}[模曲线的 Jacobian]
\label{prop:j0n-from-s2}
对紧化模曲线 $X_0(N)$，有自然同构
\[
  J_0(N)=J\bigl(X_0(N)\bigr)
  \cong \Sk_2\bigl(\GammaO(N)\bigr)^*/H_1\bigl(X_0(N),\Z\bigr).
\]
特别地，$J_0(N)$ 的维数为
\[
  \dim J_0(N)=\dim \Sk_2\bigl(\GammaO(N)\bigr)=g\bigl(X_0(N)\bigr).
\]
\end{proposition}

\begin{proof}
由\cref{thm:s2-holomorphic-differentials}，
\[
  H^0\bigl(X_0(N),\Omega^1\bigr)\cong \Sk_2\bigl(\GammaO(N)\bigr).
\]
把这个同构代入 \cref{def:jacobian} 的定义，即得
\[
  J\bigl(X_0(N)\bigr)
  =H^0\bigl(X_0(N),\Omega^1\bigr)^*/H_1\bigl(X_0(N),\Z\bigr)
  \cong \Sk_2\bigl(\GammaO(N)\bigr)^*/H_1\bigl(X_0(N),\Z\bigr).
\]
维数结论则由\cref{cor:dim-s2-genus} 立刻得到。
\end{proof}

\begin{example}[$J_0(11)$ 是一条椭圆曲线]
\label{ex:j011-elliptic}
由\cref{ex:genus-small-levels,cor:dim-s2-genus}，
\[
  g\bigl(X_0(11)\bigr)=1,
  \qquad
  \dim \Sk_2\bigl(\GammaO(11)\bigr)=1.
\]
于是根据\cref{prop:j0n-from-s2}，$J_0(11)$ 是一维 Abel 簇。
再由\cref{prop:dim1-torus-abelian,thm:genus1-jacobian}，
一维 Abel 簇就是椭圆曲线，所以 $J_0(11)$ 本身就是一条椭圆曲线。
这正是最早出现的模椭圆曲线例子之一。
\end{example}

\textbf{关键等式。}
真正值得记住的是
\[
  J_0(N)=\Sk_2\bigl(\GammaO(N)\bigr)^*/H_1\bigl(X_0(N),\Z\bigr).
\]
从这一刻开始，权 $2$ 模形式不再只是“函数空间中的向量”，
而是 Jacobian 的余切空间中的方向；
这正是 Hecke 算子能够提升为 Jacobian 自同态的根本原因。
